\documentclass[aps,preprint]{revtex4}%
\usepackage{amsfonts}%
\usepackage{amsmath}%
\setcounter{MaxMatrixCols}{30}%
\usepackage{amssymb}%
\usepackage{graphicx}
%TCIDATA{OutputFilter=latex2.dll}
%TCIDATA{Version=4.00.0.2321}
%TCIDATA{CSTFile=revtex4.cst}
%TCIDATA{Created=Tuesday, June 18, 2002 10:30:58}
%TCIDATA{LastRevised=Tuesday, June 18, 2002 10:33:25}
%TCIDATA{<META NAME="GraphicsSave" CONTENT="32">}
%TCIDATA{<META NAME="DocumentShell" CONTENT="Articles\SW\REVTeX 4">}
\newtheorem{theorem}{Theorem}
\newtheorem{acknowledgement}[theorem]{Acknowledgement}
\newtheorem{algorithm}[theorem]{Algorithm}
\newtheorem{axiom}[theorem]{Axiom}
\newtheorem{claim}[theorem]{Claim}
\newtheorem{conclusion}[theorem]{Conclusion}
\newtheorem{condition}[theorem]{Condition}
\newtheorem{conjecture}[theorem]{Conjecture}
\newtheorem{corollary}[theorem]{Corollary}
\newtheorem{criterion}[theorem]{Criterion}
\newtheorem{definition}[theorem]{Definition}
\newtheorem{example}[theorem]{Example}
\newtheorem{exercise}[theorem]{Exercise}
\newtheorem{lemma}[theorem]{Lemma}
\newtheorem{notation}[theorem]{Notation}
\newtheorem{problem}[theorem]{Problem}
\newtheorem{proposition}[theorem]{Proposition}
\newtheorem{remark}[theorem]{Remark}
\newtheorem{solution}[theorem]{Solution}
\newtheorem{summary}[theorem]{Summary}
\newenvironment{proof}[1][Proof]{\noindent\textbf{#1.} }{\ \rule{0.5em}{0.5em}}
\begin{document}
\preprint{HEP/123-qed}
\title[Short title for running header]{Long title that appears on the title page}
\author{First Author}
\affiliation{My Institution}
\author{Second Author}
\affiliation{My Institution}
\author{Third Author}
\affiliation{Other Institution}
\keywords{one two three}
\pacs{PACS number}

\begin{abstract}
Shell document for REV\TeX{} 4.

\end{abstract}
\volumeyear{year}
\volumenumber{number}
\issuenumber{number}
\eid{identifier}
\date[Date text]{date}
\received[Received text]{date}

\revised[Revised text]{date}

\accepted[Accepted text]{date}

\published[Published text]{date}

\startpage{101}
\endpage{102}
\maketitle
\tableofcontents


In order to numerically determine pure states of quantum systems one must
invariably introduce variable parameters to label states. In the molecular
context these parameters are most commonly Configuration Interaction (C.I.)
coefficients, orbital expansion coefficients and possibly exponents of the
basis functions. By imposing constraints on the variation of these
coefficients one can restrict the domain of the Schr\"{o}dinger equations
(both time dependent and independent) to sub manifolds of the Hilbert space of
the system and thus produce constrained equations whose solutions are
approximations to the full equations.

One of the most effective ways to introduce sub manifolds of a system Hilbert
space, $\mathcal{H}$, is to use a family of coherent states. The defining
property of families of coherent states is that they constitute overcomplete
continuously labelled bases \emph{which are contained} in the Hilbert space. A
class of coherent states, known as group coherent states, are generated by the
action of a group, that has an irreducible representation in $\mathcal{H}$, on
a given reference state. The orbit generated from this state forms a sub
manifold which constitutes the overcomplete basis. In these cases the
parameters can be chosen to be ones that parameterize the representation of
the coset of the group generated by the reference. A significant property of
these group coherent states is that expectation values of observables in such
states can be expressed as expectation values of "dressed" observables with
respect to the reference state. If the group acts effectively in $\mathcal{H}$
i.e. the Hilbert space is homogeneous wrt the group action, then expectation
values in any state can be developed in this manner, as the coherent state
manifold is all of the Hilbert space. In this case the Schr\"{o}dinger
equations over the coherent state manifold are the exact equations. 

Another way to generate a coherent state submanifold 

In order to obtain a sequence of approximate equations that converges to the
exact equations one can choose sequences of groups which generate a sequence
of inclusive orbits fom a given reference that converge to the full Hilbert
space, so that the limit group acts effectively. The groups that have been
used in the coherent state context have, by neccessity, been non abelian, as
their representations have to be irreducible in order to generate a resolution
of the identity.

In order to develope an approximation theory in the spirit of group coherent
state theory the essential ingredient is that the sub manifolds over which the
Schr\"{o}dinger equations are restricted to, sequentially converge to the full
Hilbert space.

$\lceil$approximations driven by:

\begin{description}
\item sequences of subgroups from a given reference ----- inclusive orbits

\item use algebras generated by sub groups on cyclic vector

\item use sequence of non cyclic vectors or cyclic vectors$\rfloor$
\end{description}

In order to consider states that do not belong to the group coherent state
manifold one needs to consider the group algebra generated by this manifold.
This leads to the generalization of the group coherent state construction to
the action of algebras and semi groups on a cyclic vector. This has the
advantage of being able to utilize abelian algebras and semi groups which have
commuting generators and thus are easier to manipulate, as it is possible to
find cyclic vectors for these algebras but not for the groups. One can combine
the coherent state concept with that of generalized bases, that are made up of
continuums of basis vectors that \emph{do not }belong to the Hilbert space, in
contrast to coherent state bases elements which are bona fide vectors
belonging to $\mathcal{H}$. A standerd example of such a bases is the one
formed by the generalized eigenvectors of the position operator. the coherIt
also has a further advantage of allowing one to consider abelian algebras and
semi groups that produce generalized bases and resolutions of the identity
from a cyclic reference state. In this article we will discuss generalized bases


\end{document}